% ===========================================================================
%  Chapitre 3 — Lire et interpréter une courbe
%  PE 2026, composante ANALYSE — RAPE, première période (« étude des graphes »)
% ===========================================================================
\bmtchapitre{Lire et interpréter une courbe}
  {Lire sur une courbe un ensemble de définition, une image, un antécédent,
   un sens de variation, un extremum, un signe et des éléments de symétrie ;
   interpréter une courbe qui représente une situation réelle.}
  {Analyse \textbullet\ environ 4 heures}

Avant d'étudier une fonction, on apprend à regarder sa courbe. Ce n'est pas
un préliminaire décoratif : à l'examen, la moitié des points de la première
question du problème se joue sur des lectures graphiques, et il est plus
facile de tracer une courbe juste quand on sait déjà lire celles des autres.

Une courbe dit beaucoup de choses d'un coup d'\oe il — où la fonction existe,
où elle monte, où elle s'annule, où elle atteint son plus haut point. Ce
chapitre consiste à mettre un mot précis sur chacune de ces observations, et
à écrire ces mots correctement.

% ---------------------------------------------------------------------------
\begin{revision}
Uniquement des acquis de seconde.

\begin{enumerate}[label=\textbf{\arabic*.}, leftmargin=*, itemsep=0.35em]
  \item Placer dans un repère les points $A(2\,;\,-1)$, $B(-3\,;\,0)$ et
        $C(0\,;\,4)$.
  \item Que signifie « le point $M$ appartient à la courbe de $f$ » ?
  \item Calculer $f(-2)$ pour $f(x) = x^{2}-3x$.
  \item Résoudre graphiquement $2x-1 = 3$.
  \item Comment reconnaît-on, sur un dessin, deux points symétriques par
        rapport à l'axe des ordonnées ?
  \item Une fonction peut-elle avoir deux images différentes pour un même
        nombre ?
\end{enumerate}
\end{revision}

\begin{automatismes}
Sans calculatrice.

\begin{enumerate}[label=\textbf{\alph*)}, leftmargin=*, itemsep=0.25em]
  \item $f(x) = 3x-5$ : calculer $f(0)$ et $f(2)$
  \item $g(x) = x^{2}$ : calculer $g(-3)$
  \item Résoudre $3x-5 = 1$
  \item $h(x) = \dfrac{1}{x}$ : $h$ est-elle définie en $0$ ?
  \item Le symétrique de $(2\,;\,5)$ par rapport à l'axe $(Oy)$
  \item Le symétrique de $(2\,;\,5)$ par rapport à l'origine
  \item Le point d'intersection de $y = 2x+6$ avec l'axe des abscisses
  \item Le point d'intersection de $y = 2x+6$ avec l'axe des ordonnées
\end{enumerate}
\end{automatismes}

\begin{corrige}[automatismes]
a) $-5$ et $1$ \quad b) $9$ \quad c) $x = 2$ \quad d) non \quad
e) $(-2\,;\,5)$ \quad f) $(-2\,;\,-5)$ \quad g) $(-3\,;\,0)$ \quad
h) $(0\,;\,6)$.
\end{corrige}

% ===========================================================================
\section{Ce qu'une courbe dit d'une fonction}

\begin{definition}[courbe représentative]
Soit $f$ une fonction définie sur une partie $\Df$ de $\R$. La
\textbf{courbe représentative} de $f$, notée $\Cf$, est l'ensemble des points
$M(x\,;\,y)$ du plan tels que
\[
  x \in \Df \quad\text{et}\quad y = f(x).
\]
Dire que le point $M(a\,;\,b)$ appartient à $\Cf$, c'est exactement dire que
$b = f(a)$.
\end{definition}

\begin{visualisation}[tout ce qui se lit sur une courbe]
\begin{center}
\begin{tikzpicture}
\begin{axis}[
  width = 12.2cm, height = 7.4cm,
  axis lines = middle, xlabel = {$x$}, ylabel = {$y$},
  xmin = -3.6, xmax = 4.6, ymin = -3.4, ymax = 4.4,
  xtick = {-3,-2,-1,1,2,3,4}, ytick = {-3,-2,-1,1,2,3,4},
  axis line style = {bmtGris},
  tick label style = {font=\scriptsize, color=bmtGris},
  grid = both, grid style = {bmtGrisClair, line width=0.3pt},
  clip = true,
]
  \addplot[bmtIndigo, line width=1.4pt, domain=-3:4, samples=200]
    {0.28*(x+2.6)*(x-0.4)*(x-3.2)};
  % maximum local
  \filldraw[bmtLaterite] (axis cs:-1.55,2.44) circle (2pt);
  \draw[bmtLaterite, dash pattern=on 2pt off 2pt]
    (axis cs:-1.55,0) -- (axis cs:-1.55,2.44) -- (axis cs:0,2.44);
  \node[bmtLaterite, font=\scriptsize, anchor=south west]
    at (axis cs:-1.5,2.5) {maximum local};
  % minimum local
  \filldraw[bmtVetiver] (axis cs:2.15,-2.42) circle (2pt);
  \draw[bmtVetiver, dash pattern=on 2pt off 2pt]
    (axis cs:2.15,0) -- (axis cs:2.15,-2.42) -- (axis cs:0,-2.42);
  \node[bmtVetiver, font=\scriptsize, anchor=north west]
    at (axis cs:2.2,-2.5) {minimum local};
  % zeros
  \filldraw[bmtRaphia] (axis cs:-2.6,0) circle (2pt);
  \filldraw[bmtRaphia] (axis cs:0.4,0) circle (2pt);
  \filldraw[bmtRaphia] (axis cs:3.2,0) circle (2pt);
  \node[bmtRaphia, font=\scriptsize, anchor=south east]
    at (axis cs:-2.55,0.12) {$f$ s'annule};
  % ordonnee a l'origine
  \filldraw[bmtIndigo] (axis cs:0,-0.93) circle (2pt);
  \node[bmtIndigo, font=\scriptsize, anchor=west]
    at (axis cs:0.15,-1.15) {$f(0)$};
\end{axis}
\end{tikzpicture}
\end{center}

Sur cette seule figure se lisent : les trois valeurs où $f$ s'annule
($-2{,}6$ ; $0{,}4$ ; $3{,}2$), l'image de $0$, les intervalles où la
fonction monte et ceux où elle descend, un maximum local atteint en
$-1{,}55$, un minimum local atteint en $2{,}15$, et le signe de $f$ — positif
là où la courbe est au-dessus de l'axe, négatif là où elle est en dessous.
Aucun calcul n'a été fait.
\end{visualisation}

\subsection{Image et antécédent}

\begin{definition}[image, antécédent]
Si $y = f(x)$, on dit que $y$ est l'\textbf{image} de $x$ par $f$, et que
$x$ est \textbf{un antécédent} de $y$.
\end{definition}

\begin{avertissement}
L'article n'est pas le même par hasard : un nombre a \textbf{une seule}
image, mais il peut avoir \textbf{plusieurs} antécédents, ou aucun.
Graphiquement : une verticale coupe la courbe en au plus un point — une
horizontale peut la couper autant de fois qu'elle veut.
\end{avertissement}

\begin{methode}{lire une image, lire un antécédent}
\begin{itemize}[leftmargin=1.3em, itemsep=0.2em]
  \item \emph{Image de $a$ :} on part de $a$ sur l'axe des abscisses, on
        monte (ou on descend) verticalement jusqu'à la courbe, on lit
        l'ordonnée.
  \item \emph{Antécédents de $b$ :} on part de $b$ sur l'axe des ordonnées,
        on se déplace horizontalement, et l'on relève \emph{toutes} les
        abscisses des points de rencontre.
\end{itemize}
\end{methode}

\subsection{Ensemble de définition et variations}

\begin{definition}[ensemble de définition lu sur une courbe]
L'\textbf{ensemble de définition} de $f$ est l'ensemble des abscisses des
points de $\Cf$ : c'est la « projection » de la courbe sur l'axe des
abscisses.
\end{definition}

\begin{definition}[sens de variation]
Sur un intervalle $I$, la fonction $f$ est \textbf{croissante} si sa courbe
monte quand on la parcourt de gauche à droite, \textbf{décroissante} si elle
descend. On résume ces informations dans un \textbf{tableau de variation}.
\end{definition}

\begin{definition}[extremum]
$f$ admet un \textbf{maximum} en $a$ sur $I$ si $f(x) \leq f(a)$ pour tout
$x$ de $I$ ; un \textbf{minimum} si $f(x) \geq f(a)$. On dit que l'extremum
est \emph{local} lorsque la propriété n'est vraie qu'au voisinage de $a$.
\end{definition}

\begin{erreurclassique}
Le maximum est une \textbf{ordonnée}, pas une abscisse. Sur la figure
ci-dessus, on écrit : « $f$ admet un maximum local égal à $2{,}44$, atteint
en $x = -1{,}55$ ». Écrire « le maximum est $-1{,}55$ » coûte le point.
\end{erreurclassique}

\subsection{Signe d'une fonction}

\begin{propriete}[signe lu sur la courbe]
$f(x) > 0$ exactement là où la courbe est \emph{strictement au-dessus} de
l'axe des abscisses ; $f(x) < 0$ là où elle est strictement en dessous ;
$f(x) = 0$ aux abscisses des points d'intersection avec cet axe.
\end{propriete}

\begin{remarque}
Résoudre graphiquement $f(x) \geq g(x)$ suit le même principe : on cherche
les abscisses où la courbe de $f$ est au-dessus de celle de $g$. Le cas
$g(x) = 0$ n'en est qu'un cas particulier, la courbe de $g$ étant alors
l'axe des abscisses lui-même.
\end{remarque}

\begin{entrainement}[lectures de base]
Toutes les questions portent sur la courbe du bloc « je visualise » de ce
chapitre, intitulé \emph{tout ce qui se lit sur une courbe}.

\exo[\niveau{1}] Donner l'ensemble de définition de $f$ tel qu'il apparaît
sur la figure.

\exo[\niveau{1}] Lire $f(0)$.

\exo[\niveau{1}] Combien $0$ a-t-il d'antécédents par $f$ ?

\exo[\niveau{2}] Donner le tableau de variation de $f$ sur
$\intervalle{-3}{4}$.

\exo[\niveau{2}] Résoudre graphiquement $f(x) \geq 0$.

\exo[\niveau{2}] Combien l'équation $f(x) = 1$ admet-elle de solutions ?

\exo[\niveau{3}] Combien l'équation $f(x) = 3$ admet-elle de solutions ?
Justifier en une phrase.
\end{entrainement}

\begin{corrige}[lectures de base]
\textbf{1.} $\Df = \intervalle{-3}{4}$ (l'intervalle effectivement tracé).
\textbf{2.} $f(0) \approx -0{,}93$.
\textbf{3.} Trois antécédents : environ $-2{,}6$ ; $0{,}4$ et $3{,}2$.
\textbf{4.} $f$ croît sur $\intervalle{-3}{-1{,}55}$, décroît sur
$\intervalle{-1{,}55}{2{,}15}$, croît sur $\intervalle{2{,}15}{4}$.
\textbf{5.} $f(x) \geq 0$ sur
$\intervalle{-2{,}6}{0{,}4}\cup\intervalle{3{,}2}{4}$.
\textbf{6.} Trois solutions : la droite $y = 1$ coupe la courbe trois fois,
puisque $1$ est compris entre le minimum local et le maximum local.
\textbf{7.} Une seule : $3$ dépasse le maximum local $2{,}44$, donc la
droite $y = 3$ ne rencontre la courbe que sur la dernière branche montante.
\end{corrige}

% ===========================================================================
\section{Parité et éléments de symétrie}

\begin{definition}[fonction paire, fonction impaire]
Soit $f$ définie sur $\Df$, où $\Df$ est symétrique par rapport à $0$.
\begin{itemize}[leftmargin=1.3em, itemsep=0.15em]
  \item $f$ est \textbf{paire} si $f(-x) = f(x)$ pour tout $x$ de $\Df$ ;
  \item $f$ est \textbf{impaire} si $f(-x) = -f(x)$ pour tout $x$ de $\Df$.
\end{itemize}
\end{definition}

\begin{theoreme}[traduction graphique]
$f$ est paire \ssi\ sa courbe admet l'axe des ordonnées pour \textbf{axe de
symétrie}. $f$ est impaire \ssi\ sa courbe admet l'origine pour
\textbf{centre de symétrie}.
\end{theoreme}

\begin{visualisation}[une paire, une impaire]
\begin{center}
\begin{tikzpicture}
\begin{axis}[
  name = P, width = 6.4cm, height = 4.9cm, axis lines = middle,
  xmin = -2.6, xmax = 2.6, ymin = -2.6, ymax = 2.9,
  xtick = \empty, ytick = \empty, axis line style = {bmtGris},
  title = {\footnotesize $f$ paire : axe $(Oy)$},
  title style = {color=bmtGris},
]
  \addplot[bmtIndigo, line width=1.3pt, domain=-2.3:2.3, samples=120]
    {0.55*x^2-0.9};
  \draw[bmtLaterite, dash pattern=on 2.5pt off 2pt, line width=0.8pt]
    (axis cs:0,-2.5) -- (axis cs:0,2.8);
  \filldraw[bmtVetiver] (axis cs:1.6,0.508) circle (1.8pt);
  \filldraw[bmtVetiver] (axis cs:-1.6,0.508) circle (1.8pt);
\end{axis}
\begin{axis}[
  name = I, at = {($(P.east)+(1.2cm,0)$)}, anchor = west,
  width = 6.4cm, height = 4.9cm, axis lines = middle,
  xmin = -2.6, xmax = 2.6, ymin = -2.6, ymax = 2.9,
  xtick = \empty, ytick = \empty, axis line style = {bmtGris},
  title = {\footnotesize $g$ impaire : centre $O$},
  title style = {color=bmtGris},
]
  \addplot[bmtIndigo, line width=1.3pt, domain=-2.1:2.1, samples=160]
    {0.28*x^3-0.15*x};
  \filldraw[bmtVetiver] (axis cs:1.8,1.363) circle (1.8pt);
  \filldraw[bmtVetiver] (axis cs:-1.8,-1.363) circle (1.8pt);
  \draw[bmtLaterite, dash pattern=on 2.5pt off 2pt, line width=0.8pt]
    (axis cs:-1.8,-1.363) -- (axis cs:1.8,1.363);
  \filldraw[bmtLaterite] (axis cs:0,0) circle (1.8pt);
\end{axis}
\end{tikzpicture}
\end{center}

À gauche, deux points de même ordonnée se font face de part et d'autre de
l'axe : c'est la parité. À droite, deux points d'ordonnées opposées sont
alignés avec l'origine, qui en est le milieu : c'est l'imparité. Reconnaître
l'une ou l'autre divise par deux le travail de tracé — il suffit d'étudier la
fonction sur les $x$ positifs, puis de compléter par symétrie.
\end{visualisation}

\begin{methode}{montrer qu'une fonction est paire ou impaire}
\begin{enumerate}[label=\textbf{\arabic*.}, leftmargin=*, itemsep=0.15em]
  \item Vérifier d'abord que $\Df$ est symétrique par rapport à $0$ : si
        $\Df = \intervalleo{0}{+\infty}$, la question ne se pose même pas.
  \item Calculer $f(-x)$ en remplaçant chaque $x$ par $-x$, sans rien
        simplifier trop vite.
  \item Comparer à $f(x)$ et à $-f(x)$.
  \item Conclure — ou conclure que la fonction n'est ni paire ni impaire,
        ce qui est le cas le plus fréquent.
\end{enumerate}
\end{methode}

\begin{exemple}[trois fonctions, trois réponses]
\textbf{a)} $f(x) = x^{2}-3$ sur $\R$ : $f(-x) = (-x)^{2}-3 = x^{2}-3
= f(x)$. La fonction est paire.

\textbf{b)} $g(x) = \dfrac{x}{x^{2}+1}$ sur $\R$ :
$g(-x) = \dfrac{-x}{x^{2}+1} = -g(x)$. La fonction est impaire.

\textbf{c)} $h(x) = x^{2}+x$ sur $\R$ : $h(-x) = x^{2}-x$, qui n'est ni
$h(x)$ ni $-h(x)$ — il suffit de le voir en $x = 1$ : $h(1) = 2$ et
$h(-1) = 0$. La fonction n'est ni paire ni impaire.
\end{exemple}

\begin{propriete}[axe et centre de symétrie ailleurs qu'à l'origine]
La droite d'équation $x = a$ est axe de symétrie de $\Cf$ \ssi, pour tout
$h$ tel que $a\pm h$ soit dans $\Df$,
\[
  f(a-h) = f(a+h).
\]
Le point $\Omega(a\,;\,b)$ est centre de symétrie de $\Cf$ \ssi
\[
  f(a-h)+f(a+h) = 2b .
\]
\end{propriete}

\begin{entrainement}[parité et symétries]
\exo[\niveau{1}] $f(x) = x^{4}-2x^{2}$ : étudier la parité.

\exo[\niveau{1}] $g(x) = x^{3}+x$ : étudier la parité.

\exo[\niveau{2}] $h(x) = \dfrac{1}{x^{2}-4}$ : préciser $\Df$, puis étudier
la parité.

\exo[\niveau{2}] $k(x) = 2x+1$ : cette fonction est-elle paire, impaire, ni
l'une ni l'autre ?

\exo[\niveau{3}] Montrer que la droite d'équation $x = 1$ est axe de
symétrie de la courbe de $f(x) = x^{2}-2x+5$.

\exo[\niveau{3}] Montrer que le point $\Omega(0\,;\,3)$ est centre de
symétrie de la courbe de $g(x) = x^{3}+3$.
\end{entrainement}

\begin{corrige}[parité et symétries]
\textbf{1.} $f(-x) = x^{4}-2x^{2} = f(x)$ : paire.
\textbf{2.} $g(-x) = -x^{3}-x = -g(x)$ : impaire.
\textbf{3.} $\Df = \R\setminus\{-2\,;\,2\}$, symétrique par rapport à $0$ ;
$h(-x) = h(x)$ : paire.
\textbf{4.} $k(-x) = -2x+1$ : ni l'un ni l'autre (par exemple $k(1) = 3$ et
$k(-1) = -1$).
\textbf{5.} $f(1-h) = (1-h)^{2}-2(1-h)+5 = h^{2}+4$ et
$f(1+h) = (1+h)^{2}-2(1+h)+5 = h^{2}+4$ : les deux sont égaux.
\textbf{6.} $g(-h)+g(h) = -h^{3}+3+h^{3}+3 = 6 = 2\times3$.
\end{corrige}

\begin{qcm}
\exo Si $f(-x) = f(x)$ pour tout $x$, la courbe de $f$ admet :
\textbf{a)} l'origine pour centre de symétrie \quad
\textbf{b)} l'axe des ordonnées pour axe de symétrie \quad
\textbf{c)} l'axe des abscisses pour axe de symétrie \quad
\textbf{d)} aucune symétrie \reponseqcm{b}

\exo Un nombre réel peut avoir, par une fonction :
\textbf{a)} plusieurs images \quad \textbf{b)} exactement une image ou
aucune \quad \textbf{c)} toujours exactement une image \quad
\textbf{d)} exactement deux images \reponseqcm{b}

\exo Sur $\intervalle{-3}{4}$, une fonction croissante puis décroissante
admet en son point de retournement :
\textbf{a)} un minimum \quad \textbf{b)} un maximum \quad
\textbf{c)} un zéro \quad \textbf{d)} rien de particulier \reponseqcm{b}

\exo La courbe de $f$ coupe l'axe des ordonnées en un point d'ordonnée :
\textbf{a)} $f(1)$ \quad \textbf{b)} $0$ \quad \textbf{c)} $f(0)$ \quad
\textbf{d)} l'antécédent de $0$ \reponseqcm{c}

\exo Si $f$ est impaire et définie en $0$, alors nécessairement :
\textbf{a)} $f(0) = 1$ \quad \textbf{b)} $f(0) = 0$ \quad
\textbf{c)} $f(0) > 0$ \quad \textbf{d)} on ne peut rien dire
\reponseqcm{b}
\end{qcm}

% ===========================================================================
\begin{lecturegraphique}[une seule figure, huit questions]
\begin{center}
\begin{tikzpicture}
\begin{axis}[
  width = 11.4cm, height = 6.6cm,
  axis lines = middle, xlabel = {$x$}, ylabel = {$y$},
  xmin = -4.4, xmax = 5.4, ymin = -3.4, ymax = 4.4,
  xtick = {-4,-3,-2,-1,1,2,3,4,5}, ytick = {-3,-2,-1,1,2,3,4},
  axis line style = {bmtGris},
  tick label style = {font=\scriptsize, color=bmtGris},
  grid = both, grid style = {bmtGrisClair, line width=0.3pt}, clip = true,
]
  \addplot[bmtIndigo, line width=1.4pt, domain=-4:5, samples=220]
    {2.6*sin(deg(0.62*x+0.5))+0.55};
  \addplot[bmtLaterite, line width=1pt, dash pattern=on 3pt off 2pt,
           domain=-4.2:5.2] {1};
\end{axis}
\end{tikzpicture}
\end{center}

La courbe ci-dessus est celle d'une fonction $g$ tracée sur
$\intervalle{-4}{5}$. La droite en pointillé a pour équation $y = 1$.

\begin{enumerate}[label=\textbf{\arabic*.}, leftmargin=*, itemsep=0.25em]
  \item Donner l'ensemble de définition de $g$ tel qu'il apparaît.
  \item Donner une valeur approchée de $g(0)$ et de $g(-2)$.
  \item Donner le nombre d'antécédents de $1$ par $g$, et une valeur
        approchée de chacun.
  \item Dresser le tableau de variation de $g$.
  \item Donner le maximum de $g$ et l'abscisse en laquelle il est atteint.
  \item Résoudre graphiquement $g(x) \geq 1$.
  \item Résoudre graphiquement $g(x) < 0$.
  \item La fonction $g$ semble-t-elle paire ? impaire ? Justifier.
\end{enumerate}
\end{lecturegraphique}

\begin{corrige}[lecture graphique]
\textbf{1.} $\Df = \intervalle{-4}{5}$.
\textbf{2.} $g(0) \approx 1{,}8$ et $g(-2) \approx -1{,}2$.
\textbf{3.} Trois antécédents, d'abscisses approximatives $-2{,}5$ ;
$-0{,}6$ et $4{,}5$.
\textbf{4.} $g$ décroît sur $\intervalle{-4}{-3{,}3}$, croît sur
$\intervalle{-3{,}3}{1{,}7}$, décroît sur $\intervalle{1{,}7}{5}$
(valeurs lues, donc approchées).
\textbf{5.} Le maximum vaut environ $3{,}15$, atteint pour
$x \approx 1{,}7$.
\textbf{6.} $g(x) \geq 1$ sur $\intervalle{-0{,}6}{4{,}5}$ environ.
\textbf{7.} $g(x) < 0$ sur $\intervalleof{-4}{-1{,}5}$ environ, ainsi que
sur la fin de l'intervalle si la courbe repasse sous l'axe — ici elle ne le
fait pas avant $x = 5$.
\textbf{8.} Ni l'une ni l'autre : $g(2) \neq g(-2)$, donc elle n'est pas
paire ; et $g(0) \neq 0$, donc elle n'est pas impaire.
\end{corrige}

% ===========================================================================
\begin{situation}[la courbe de température d'un séchoir à vanille]
Un séchoir solaire est allumé à $6$ heures du matin. La courbe ci-dessous
donne la température $T$, en degrés Celsius, en fonction de l'heure $t$.

\begin{center}
\begin{tikzpicture}
\begin{axis}[
  width = 10.6cm, height = 5.4cm,
  axis lines = left, xlabel = {$t$ (heures)}, ylabel = {$T$ (°C)},
  xmin = 5.5, xmax = 19, ymin = 15, ymax = 60,
  xtick = {6,8,10,12,14,16,18}, ytick = {20,30,40,50},
  axis line style = {bmtGris},
  tick label style = {font=\scriptsize, color=bmtGris},
  grid = both, grid style = {bmtGrisClair, line width=0.3pt}, clip = true,
]
  \addplot[bmtRaphia, line width=1.4pt, domain=6:18, samples=140]
    {22+30*sin(deg(3.1416*(x-6)/13.5))};
  \addplot[bmtIndigo, line width=1pt, dash pattern=on 3pt off 2pt,
           domain=5.5:19] {45};
\end{axis}
\end{tikzpicture}
\end{center}

Le séchage de la vanille exige une température d'au moins $45$ °C.

\begin{enumerate}[label=\textbf{\arabic*.}, leftmargin=*, itemsep=0.3em]
  \item Quelle est la température à l'allumage ? à midi ?
  \item À quelle heure la température est-elle maximale, et que vaut ce
        maximum ?
  \item Sur quel intervalle de temps la température est-elle supérieure ou
        égale à $45$ °C ? Combien d'heures de séchage utile cela
        représente-t-il ?
  \item La fonction $T$ est-elle croissante sur tout l'intervalle
        $\intervalle{6}{18}$ ? Répondre par une phrase, en distinguant les
        périodes.
\end{enumerate}
\end{situation}

\begin{corrige}[la courbe de température d'un séchoir à vanille]
\textbf{1.} À $6$ h, $T = 22$ °C ; à midi, $T \approx 50$ °C.

\textbf{2.} Le maximum, environ $52$ °C, est atteint vers $12$ h $45$.

\textbf{3.} La courbe est au-dessus de la droite $y = 45$ entre environ
$10$ h et $15$ h $30$, soit à peu près cinq heures et demie de séchage
utile.

\textbf{4.} Non : la température croît de $6$ h jusqu'à environ $12$ h $45$,
puis décroît jusqu'à $18$ h. On ne peut donc pas dire que $T$ est croissante
sur $\intervalle{6}{18}$ ; elle l'est seulement sur
$\intervalle{6}{12{,}75}$.
\end{corrige}

% ===========================================================================
\begin{annales}
La lecture graphique n'est jamais un exercice séparé à l'examen : elle est
la première ou la dernière question du problème. Les extraits ci-dessous sont
ceux qui ne demandent que de lire ou d'interpréter, sans aucune notion
postérieure à ce chapitre.

\exoann{Baccalauréat, Madagascar, série A, session 2016 — problème,
question 4 (extrait)}
On donne le tableau de valeurs suivant, arrondi à $10^{-2}$ près, d'une
fonction $f$ définie sur $\intervalleo{0}{+\infty}$ :

\begin{center}\footnotesize
\begin{tabular}{|c|c|c|c|c|}
\hline
$x$ & $2$ & $2{,}72$ & $4$ & $6$\\
\hline
$f(x)$ & $1{,}31$ & $1{,}72$ & $2{,}61$ & $4{,}21$\\
\hline
\end{tabular}
\end{center}

Placer les quatre points correspondants dans un repère d'unité $1$ cm, puis
dire, à la seule lecture de ces points, si la fonction semble croissante ou
décroissante sur $\intervalle{2}{6}$.

\exoann{Baccalauréat, Madagascar, série A, session 2013 — problème,
question 3.a) (extrait)}
Une fonction $f$ vérifie $f(0) = 2$, $f(1) \approx 2{,}92$ et
$f(2) \approx 3{,}52$, et sa courbe admet la droite d'équation $y = 4$ pour
asymptote horizontale. Placer ces trois points, tracer la droite, et dire ce
que la courbe ne peut pas faire.

\exoann{Baccalauréat, Madagascar, série A, session 2017 — problème,
question 4.a) (extrait)}
On admet que le point $I\left(0\,;\,\frac52\right)$ est centre de symétrie de
la courbe d'une fonction $f$. Sachant que $f(-1) \approx 2{,}27$, en déduire
$f(1)$ sans calcul.
\end{annales}

\begin{corrige}[annales]
\textbf{1.} Les quatre ordonnées augmentent quand $x$ augmente : la fonction
semble croissante sur $\intervalle{2}{6}$. On ne peut rien affirmer de plus
à partir d'un tableau de valeurs — il faudra le tableau de variation pour le
démontrer.

\textbf{2.} Les trois points montent, et la droite $y = 4$ est asymptote :
la courbe s'en approche sans jamais l'atteindre. Elle ne peut donc pas
dépasser l'ordonnée $4$ ni la toucher.

\textbf{3.} Si $I(0\,;\,b)$ est centre de symétrie, alors
$f(-1)+f(1) = 2b = 5$. Donc $f(1) = 5-2{,}27 = 2{,}73$.
\end{corrige}

% ===========================================================================
\begin{vraifaux}
\exo Si une courbe coupe l'axe des abscisses en trois points, l'équation
$f(x) = 0$ a trois solutions.

\exo Une fonction croissante sur $\intervalle{0}{2}$ et croissante sur
$\intervalle{2}{5}$ est croissante sur $\intervalle{0}{5}$.

\exo Si $f$ est paire, il suffit d'étudier $f$ sur les réels positifs.

\exo Le maximum d'une fonction est atteint au point le plus à droite de sa
courbe.

\exo Une horizontale peut couper la courbe d'une fonction en deux points.
\end{vraifaux}

\begin{corrige}[vrai ou faux]
\textbf{1.} Vrai — chaque point d'intersection donne une solution, et
réciproquement.
\textbf{2.} Vrai ici, parce que les deux intervalles se recollent en $2$ où
la fonction est définie ; attention, ce recollement échoue si la fonction
n'est pas définie en $2$.
\textbf{3.} Vrai — la partie négative se déduit par symétrie d'axe $(Oy)$.
\textbf{4.} Faux — il est atteint au point le plus \emph{haut}.
\textbf{5.} Vrai — c'est la verticale, et elle seule, qui ne peut couper la
courbe qu'une fois.
\end{corrige}

% ===========================================================================
\begin{jemeteste}
Vingt minutes, la figure « une seule figure, huit questions » sous les yeux.

\begin{enumerate}[label=\textbf{\arabic*.}, leftmargin=*, itemsep=0.3em]
  \item Écrire ce que signifie, sans le mot « courbe », l'égalité
        $g(3) = 2$.
  \item Donner un intervalle sur lequel $g$ est décroissante.
  \item Étudier la parité de $f(x) = \dfrac{x^{2}+1}{x^{2}-1}$.
  \item Montrer que le point $\Omega(2\,;\,0)$ est centre de symétrie de la
        courbe de $h(x) = (x-2)^{3}$.
  \item Une fonction admet un maximum égal à $7$ atteint en $x = -1$.
        Traduire cette phrase par une inégalité valable pour tout $x$ de
        l'ensemble de définition.
\end{enumerate}
\end{jemeteste}

\begin{corrige}[je me teste]
\textbf{1.} L'image de $3$ par $g$ vaut $2$ ; autrement dit, $3$ est un
antécédent de $2$.
\textbf{2.} Par exemple $\intervalle{2}{5}$.
\textbf{3.} $\Df = \R\setminus\{-1\,;\,1\}$, symétrique, et
$f(-x) = f(x)$ : la fonction est paire.
\textbf{4.} $h(2-t)+h(2+t) = -t^{3}+t^{3} = 0 = 2\times0$.
\textbf{5.} $f(x) \leq 7$ pour tout $x$, avec égalité pour $x = -1$.
\end{corrige}

% ===========================================================================
\begin{commeaubac}
\bmtsource{Exercice original au format de l'épreuve — série A, options A1 et
A2 \textbullet\ 5 points}
La courbe $\Cf$ ci-dessous est celle d'une fonction $f$ définie sur
$\intervalle{-3}{4}$.

\begin{center}
\begin{tikzpicture}
\begin{axis}[
  width = 10.2cm, height = 6cm,
  axis lines = middle, xlabel = {$x$}, ylabel = {$y$},
  xmin = -3.6, xmax = 4.6, ymin = -4.4, ymax = 5.4,
  xtick = {-3,-2,-1,1,2,3,4}, ytick = {-4,-3,-2,-1,1,2,3,4,5},
  axis line style = {bmtGris},
  tick label style = {font=\scriptsize, color=bmtGris},
  grid = both, grid style = {bmtGrisClair, line width=0.3pt}, clip = true,
]
  \addplot[bmtIndigo, line width=1.4pt, domain=-3:4, samples=200]
    {0.5*x^3-1.5*x^2+2};
\end{axis}
\end{tikzpicture}
\end{center}

\begin{enumerate}[label=\textbf{\arabic*.}, leftmargin=*, itemsep=0.25em]
  \item Donner l'ensemble de définition de $f$ et lire $f(0)$.
        \bareme{0,5 pt}
  \item Dresser le tableau de variation de $f$ sur $\intervalle{-3}{4}$.
        \bareme{1,5 pt}
  \item Donner le maximum local et le minimum local de $f$, en précisant en
        quelles abscisses ils sont atteints. \bareme{1 pt}
  \item Déterminer graphiquement le nombre de solutions de l'équation
        $f(x) = 0$. \bareme{1 pt}
  \item Résoudre graphiquement l'inéquation $f(x) \leq 0$. \bareme{1 pt}
\end{enumerate}
\end{commeaubac}

\begin{corrige}[comme au baccalauréat]
\textbf{1.} $\Df = \intervalle{-3}{4}$ et $f(0) = 2$.

\textbf{2.} La courbe monte de $x = -3$ à $x = 0$, descend de $0$ à $2$,
remonte de $2$ à $4$ :

\begin{center}\footnotesize
\begin{tabular}{|c|ccccccc|}
\hline
$x$ & $-3$ & & $0$ & & $2$ & & $4$\\
\hline
$f$ & $-11{,}5$ & $\nearrow$ & $2$ & $\searrow$ & $-2$ & $\nearrow$ & $2$\\
\hline
\end{tabular}
\end{center}

\textbf{3.} Maximum local $2$ atteint en $x = 0$ ; minimum local $-2$
atteint en $x = 2$.

\textbf{4.} La courbe coupe l'axe des abscisses en trois points : l'équation
$f(x) = 0$ a trois solutions, d'abscisses approximatives $-1{,}1$ ; $1{,}3$
et $3{,}8$.

\textbf{5.} $f(x) \leq 0$ là où la courbe est sur l'axe ou en dessous :
$\mathcal S \approx \intervalle{-3}{-1{,}1}\cup\intervalle{1{,}3}{3{,}8}$.
\end{corrige}

\begin{notepourlaclasse}
Quatre heures suffisent, et il ne faut pas les dépasser : ce chapitre est
sans difficulté technique, mais il conditionne la note du problème. Le temps
utile est celui passé à faire \emph{écrire} les réponses, pas à faire lire
les figures : les élèves lisent juste et rédigent faux.

Trois formulations à exiger dès ce chapitre, et à ne plus jamais laisser
passer ensuite : « $f$ est croissante \emph{sur} l'intervalle… » (jamais
« la courbe est croissante »), « le maximum \emph{vaut} $2$, \emph{atteint
en} $x = 0$ », et « $f(x) \geq 0$ \emph{sur} … » avec un intervalle, jamais
une inégalité recopiée.

Le bloc « lecture graphique » se prête bien à un travail en binôme avec une
seule feuille pour deux : celui qui tient le crayon n'est pas celui qui lit.
\end{notepourlaclasse}
